\begin{filecontents*}{discretization-references.bib}
@Article{FarmerToda2017QE,
  author  = {Farmer, Leland E. and Toda, Alexis Akira},
  title   = {Discretizing Nonlinear, Non-{G}aussian {M}arkov Processes with Exact Conditional Moments},
  journal = {Quantitative Economics},
  year    = {2017},
  volume  = {8},
  number  = {2},
  pages   = {651--683},
  month   = jul,
  doi     = {10.3982/QE737},
}

@Article{TanakaToda2013EL,
  author  = {Tanaka, Ken'ichiro and Toda, Alexis Akira},
  title   = {Discrete Approximations of Continuous Distributions by Maximum Entropy},
  journal = {Economics Letters},
  year    = {2013},
  volume  = {118},
  number  = {3},
  pages   = {445--450},
  month   = mar,
  doi     = {10.1016/j.econlet.2012.12.020},
}

@Article{TanakaToda2015SINUM,
  author  = {Tanaka, Ken'ichiro and Toda, Alexis Akira},
  title   = {Discretizing Distributions with Exact Moments: Error Estimate and Convergence Analysis},
  journal = {{SIAM} Journal on Numerical Analysis},
  year    = {2015},
  volume  = {53},
  number  = {5},
  pages   = {2158--2177},
  month   = jan,
  doi     = {10.1137/140971269},
}

@Article{Toda2021CE,
  author  = {Toda, Alexis Akira},
  title   = {Data-based Automatic Discretization of Nonparametric Distributions},
  journal = {Computational Economics},
  year    = {2021},
  volume  = {57},
  pages   = {1217--1235},
  month   = apr,
  doi     = {10.1007/s10614-020-10012-6},
}
\end{filecontents*}

\documentclass[a4paper,12pt]{article}
\pdfoutput=1
\usepackage{amsmath,amssymb}

\usepackage[T1]{fontenc}
\usepackage{lmodern}
\usepackage{xcolor}
\usepackage{booktabs}
\usepackage{tabularx}
\usepackage{array}
\usepackage{listings}
\usepackage[inline,shortlabels]{enumitem}
\usepackage[margin=1.15in]{geometry}
\usepackage{setspace}
\usepackage{microtype}
\usepackage{csquotes}
\usepackage{hyperref}
\hypersetup{
  colorlinks=true,
  linkcolor=blue!55!black,
  citecolor=blue!55!black,
  urlcolor=magenta!70!black,
  pdftitle={Farmer-Toda Discretization: MATLAB and Python User Guide},
  pdfauthor={Alexis Akira Toda}
}
\usepackage{doi}

\usepackage[
  style=ext-authoryear-comp,
  sorting=nyt,
  dashed=false,
  maxcitenames=2,
  maxbibnames=99,
  uniquelist=false,
  uniquename=false,
  giveninits=true,
  natbib,
  date=year
]{biblatex}

\AtBeginRefsection{\GenRefcontextData{sorting=ynt}}
\AtEveryCite{\localrefcontext[sorting=ynt]}
\DeclareFieldFormat{pages}{#1}
\renewbibmacro{in:}{%
  \ifentrytype{article}{}{\printtext{\bibstring{in}\intitlepunct}}}
\DeclareFieldFormat[
  article,inbook,incollection,inproceedings,patent,thesis,unpublished
]{titlecase:title}{\MakeSentenceCase*{#1}}
\addbibresource{discretization-references.bib}

% Minimal mathematical definitions used below.
\newcommand{\e}{\mathrm{e}}
\newcommand{\diff}{\mathop{}\!\mathrm{d}}
\newcommand{\R}{\mathbb{R}}
\newcommand{\CE}[1]{\operatorname{E}\!\left[#1\right]}
\newcommand{\Cvar}[1]{\operatorname{Var}\!\left[#1\right]}

\setlist[enumerate,1]{label=(\roman*)}
\setlist[itemize]{itemsep=0.25em,topsep=0.4em}
\onehalfspacing
\setlength{\emergencystretch}{2em}

\definecolor{codebackground}{RGB}{247,248,250}
\definecolor{codeframe}{RGB}{210,214,220}
\lstdefinestyle{guidecode}{
  basicstyle=\ttfamily\small,
  backgroundcolor=\color{codebackground},
  frame=single,
  rulecolor=\color{codeframe},
  breaklines=true,
  breakatwhitespace=false,
  columns=fullflexible,
  keepspaces=true,
  showstringspaces=false,
  xleftmargin=0.35em,
  xrightmargin=0.35em,
  aboveskip=0.8em,
  belowskip=0.8em
}
\lstset{style=guidecode}

\newcommand{\pkg}{\texttt{discretization}}
\newcommand{\matlab}{\textsc{Matlab}}
\newcommand{\python}{\textsc{Python}}
\newcommand{\code}[1]{\texttt{#1}}
\newcommand{\matgmar}{\texttt{discretization.\allowbreak
discreteGaussian\allowbreak MixtureAR}}
\newcommand{\matsv}{\texttt{discretization.\allowbreak
discreteStochastic\allowbreak VolatilityAR}}
\newcommand{\matdist}{\texttt{discretization.\allowbreak
momentMatched\allowbreak Distribution}}
\newcommand{\pygmar}{\texttt{discrete\_gaussian\_mixture\_\allowbreak ar}}
\newcommand{\pysv}{\texttt{discrete\_stochastic\_\allowbreak volatility\_ar}}
\newcommand{\pydist}{\texttt{moment\_matched\_\allowbreak distribution}}
\newcommand{\matmaxent}{\texttt{discretization.\allowbreak
maximumEntropy\allowbreak Weights}}
\newcommand{\pymaxent}{\texttt{maximum\_entropy\_\allowbreak weights}}
\newcommand{\matgmq}{\texttt{discretization.\allowbreak
gaussianMixture\allowbreak Quadrature}}
\newcommand{\pygmq}{\texttt{gaussian\_mixture\_\allowbreak quadrature}}
\newcommand{\matnpgq}{\texttt{discretization.\allowbreak
dataDriven\allowbreak GaussianQuadrature}}
\newcommand{\pynpgq}{\texttt{data\_driven\_gaussian\_\allowbreak quadrature}}
\newcommand{\matcirpdf}{\texttt{discretization.\allowbreak
cirTransition\allowbreak Density}}
\newcommand{\pycirpdf}{\texttt{cir\_transition\_\allowbreak density}}

\title{Farmer--Toda Discretization\\
\large MATLAB and Python User Guide}
\author{Alexis Akira Toda\thanks{Department of Economics, Emory University.
Email: \href{mailto:alexis.akira.toda@emory.edu}{alexis.akira.toda@emory.edu}.}}
\date{First version: June 2, 2019\\Updated: \today}

\begin{document}
\maketitle

\begin{abstract}
This guide documents the Farmer--Toda routines for approximating continuous
distributions and nonlinear Markov processes by finite-state models that match
selected moments. The repository provides a namespaced \matlab{} package and a
\python{} implementation with common model coverage, automated tests, and
cross-language numerical reference checks. The guide describes installation,
the public application programming interface (API), model-specific inputs and
outputs, grid and moment choices, numerical safeguards, compatibility names,
and the construction of new discretization routines.

\medskip
\noindent\textbf{Keywords:} discretization; Markov chain approximation;
maximum entropy; quadrature; moment matching
\end{abstract}

\tableofcontents

\section{Introduction}

The \pkg{} repository implements the moment-preserving discretization method
developed by \textcite{FarmerToda2017QE}, building on the maximum-entropy
construction of \textcite{TanakaToda2013EL} and the convergence analysis in
\textcite{TanakaToda2015SINUM}. The central idea is to start from a finite grid
and a prior distribution, then adjust the probability weights so that selected
moments agree with their targets.

The modernized repository has two user-facing implementations:
\begin{itemize}
  \item a \matlab{} package called \code{discretization}, invoked through
  names such as \code{discretization.discreteAR}; and
  \item a \python{} package with snake-case names such as
  \code{discrete\_ar}.
\end{itemize}
Historical \matlab{} entry points remain available as compatibility wrappers,
but new code should use the package-qualified names in this guide. The two
implementations are tested against deterministic \matlab{} R2024b reference
outputs. Numerical solvers can produce small platform-dependent differences,
so parity is assessed with explicit tolerances rather than bit-for-bit
identity.

Whenever these routines are used in research, please cite at least
\textcite{FarmerToda2017QE}; citations to \textcite{TanakaToda2013EL} and
\textcite{TanakaToda2015SINUM} are also appropriate when the
maximum-entropy construction is material to the application.

\section{Getting started}
\label{sec:getting-started}

\subsection{MATLAB}

Add only the repository root to the \matlab{} path. Package and private helper
resolution then occur automatically.

\begin{lstlisting}[language=Matlab]
addpath('path/to/discretization-master')

[P, X] = discretization.discreteAR(0, 0.9, 0.1, 9);
\end{lstlisting}

The implementation uses optimization and probability-distribution routines,
including \code{fminunc}, \code{fmincon}, \code{lsqnonneg}, \code{normpdf},
and \code{norminv}. A standard installation therefore requires the
Optimization Toolbox and the Statistics and Machine Learning Toolbox.

\subsection{Python}

Install the package from the repository's \code{python} directory. Python
3.10 or newer, NumPy, and SciPy are required.

\begin{lstlisting}
python -m pip install ./python
\end{lstlisting}

\begin{lstlisting}[language=Python]
from discretization import discrete_ar

P, grid = discrete_ar(
    mean=0.0,
    persistence=0.9,
    innovation_std=0.1,
    state_count=9,
)
\end{lstlisting}

\subsection{Output convention}

All process routines return a row-stochastic transition matrix $P$. Entry
$P_{ij}$ is the probability of moving from current state $i$ to destination
state $j$. A univariate \python{} grid is one-dimensional. Multivariate and
joint-state grids use one row per state variable and one column per Markov
state. In \matlab{}, process grids follow the same row-by-variable convention.

Tensor-product states are ordered with the last state dimension changing
fastest. This detail matters when a transition matrix is combined with a value
function, simulation code, or another tensor-product operator.

\section{Public API}
\label{sec:api}

Table \ref{tab:api} lists the main public functions. The long descriptive names
are intentional: they distinguish an autoregression with Gaussian-mixture
innovations from a mixture autoregression, and they identify the particular
stochastic-volatility model being approximated.

\begin{table}[htbp]
\centering
\caption{Canonical public API}
\label{tab:api}
\footnotesize
\begin{tabularx}{\textwidth}{@{}>{\raggedright\arraybackslash}p{0.24\textwidth}
>{\raggedright\arraybackslash}X
>{\raggedright\arraybackslash}X@{}}
\toprule
Purpose & \matlab{} & \python{} \\
\midrule
Gaussian AR(1) & \code{discretization.discreteAR}
& \code{discrete\_ar} \\
Gaussian VAR(1) & \code{discretization.discreteVAR}
& \code{discrete\_var} \\
Cox--Ingersoll--Ross & \code{discretization.discreteCIR}
& \code{discrete\_cir} \\
AR with mixture shocks & \matgmar & \pygmar \\
AR with stochastic volatility & \matsv & \pysv \\
Distribution from moments & \matdist & \pydist \\
Maximum-entropy weights & \matmaxent & \pymaxent \\
Gaussian-mixture quadrature &
\matgmq & \pygmq \\
Data-driven quadrature &
\matnpgq & \pynpgq \\
CIR transition density & \matcirpdf & \pycirpdf \\
\bottomrule
\end{tabularx}
\end{table}

\section{Maximum-entropy moment matching}
\label{sec:maxent}

Let $D=(x_1,\ldots,x_N)$ be a fixed grid, let $q_n\geq 0$ be prior weights,
and let $T(x_n)\in\R^L$ be the moment functions. The target moment vector is
$\overline T\in\R^L$. The routines minimize Kullback--Leibler divergence from
the prior subject to the probability and moment restrictions. The dual
representation gives
\begin{equation}
  p_n(\lambda)
  = \frac{q_n\exp\{\lambda^{\mathsf T}[T(x_n)-\overline T]\}}
  {\sum_{j=1}^N q_j
  \exp\{\lambda^{\mathsf T}[T(x_j)-\overline T]\}}.
  \label{eq:maxent-probability}
\end{equation}
The gradient of the log normalizer is the moment error. The Hessian is the
weighted covariance matrix of the moment functions.

The modernized implementation evaluates the dual objective with a log-sum-exp
shift, supplies analytic first and second derivatives, caches solver settings,
and retains the best finite candidate across solver attempts. Before expensive
higher-order solves, selected routines test whether the target lies in the
convex hull generated by the finite grid. An infeasible target cannot be
matched by any probability vector on that grid.

The low-level interfaces are
\begin{lstlisting}[language=Matlab]
[p, lambda, momentError] = ...
    discretization.maximumEntropyWeights(D, T, TBar, q, lambda0);
\end{lstlisting}
and
\begin{lstlisting}[language=Python]
p, dual, moment_error = maximum_entropy_weights(
    grid, moment_function, target_moments, prior, dual0
)
\end{lstlisting}
where \code{q}/\code{prior} and \code{lambda0}/\code{dual0} are optional.
Grid points occupy columns in \matlab{} and the last array axis in \python{}.

\section{Process discretization routines}
\label{sec:processes}

\subsection{Gaussian AR(1)}

The Gaussian autoregression is
\begin{equation}
  x_t=(1-\rho)\mu+\rho x_{t-1}+\varepsilon_t,
  \qquad \varepsilon_t\sim N(0,\sigma^2).
\end{equation}
The canonical calls are
\begin{lstlisting}[language=Matlab]
[P, X] = discretization.discreteAR( ...
    mu, rho, sigma, Nm, method, nMoments, nSigmas);
\end{lstlisting}
\begin{lstlisting}[language=Python]
P, grid = discrete_ar(
    mean, persistence, innovation_std, state_count,
    method="even", moment_count=2, grid_width=None,
)
\end{lstlisting}
The default is an evenly spaced grid and two conditional moments. The
\matlab{} routine also exposes its quadrature-grid variants. The current
\python{} \code{discrete\_ar} interface supports the evenly spaced grid.

\subsection{Gaussian VAR(1)}

For
\begin{equation}
  y_t=b+B y_{t-1}+\eta_t,
  \qquad \eta_t\sim N(0,\Psi),
\end{equation}
use
\begin{lstlisting}[language=Matlab]
[P, X] = discretization.discreteVAR( ...
    b, B, Psi, Nm, nMoments, method, nSigmas);
\end{lstlisting}
\begin{lstlisting}[language=Python]
P, grid = discrete_var(
    constant, lag, innovation_covariance, state_count,
    moment_count=2, method="even", grid_width=None,
)
\end{lstlisting}
Here \code{Nm}/\code{state\_count} is the number of grid points in each
dimension, so a $d$-dimensional VAR has $N_m^d$ states and an
$N_m^d\times N_m^d$ transition matrix. The covariance matrix $\Psi$ must be
positive definite and the spectral radius of $B$ must be below one.

Available methods are \code{'even'}, \code{'quadrature'} (Gauss--Hermite),
and \code{'quantile'}. The quantile method is retained for replication but is
generally poor. For a weakly persistent process, quadrature can be accurate;
for persistent processes, the evenly spaced grid is usually more reliable.
The implementation rotates standardized innovations to make unconditional
component variances as similar as possible. The two-dimensional rotation has a
closed-form solution.

\subsection[AR(p) with Gaussian-mixture shocks]{AR($p$) with Gaussian-mixture shocks}

Consider
\begin{equation}
  x_t-\mu=\sum_{j=1}^p a_j(x_{t-j}-\mu)+\varepsilon_t,
\end{equation}
where $\varepsilon_t$ is an i.i.d. Gaussian mixture. The interfaces are
\begin{lstlisting}[language=Matlab]
[P, X] = discretization.discreteGaussianMixtureAR( ...
    mu, A, pC, muC, sigmaC, Nm, nMoments, method, nSigmas);
\end{lstlisting}
\begin{lstlisting}[language=Python]
P, grid = discrete_gaussian_mixture_ar(
    mean, ar_coefficients, mixture_probabilities,
    mixture_means, mixture_standard_deviations, state_count,
    moment_count=2, method="even", grid_width=None,
)
\end{lstlisting}
The vectors \code{pC}, \code{muC}, and \code{sigmaC} contain component
probabilities, means, and standard deviations. Probabilities must be
nonnegative and sum to one. The AR companion matrix must be stable.

Grid methods are evenly spaced, Gauss--Legendre, Clenshaw--Curtis,
Gauss--Hermite, and Gaussian-mixture quadrature. \matlab{} method strings use
hyphens and \code{'GMQ'}; \python{} uses
\code{'gauss\_legendre'}, \code{'clenshaw\_curtis'},
\code{'gauss\_hermite'}, and \code{'gaussian\_mixture'}, while accepting the
corresponding hyphenated aliases. Unless the process is weakly persistent and
a quadrature prior is especially useful, the evenly spaced grid is the safest
default. The number of states grows as $N_m^p$, so small $p$ is essential.

\subsection{AR(1) with stochastic volatility}

The joint model is
\begin{align}
  y_t &= \lambda y_{t-1}+u_t,
  &u_t&\sim N(0,\e^{x_t}),\\
  x_t &= (1-\rho)\mu+\rho x_{t-1}+\epsilon_t,
  &\epsilon_t&\sim N(0,\sigma_E^2).
\end{align}
Use
\begin{lstlisting}[language=Matlab]
[P, yxGrid] = discretization.discreteStochasticVolatilityAR( ...
    lambda, rho, sigmaU, sigmaE, Ny, Nx, method, nSigmaY);
\end{lstlisting}
\begin{lstlisting}[language=Python]
P, joint_grid = discrete_stochastic_volatility_ar(
    persistence, volatility_persistence,
    unconditional_innovation_std, volatility_innovation_std,
    state_count, volatility_state_count,
    method="even", grid_width=None,
)
\end{lstlisting}
The returned grid has the $y$ state first and log variance second. There are
$N_yN_x$ joint states. The argument \code{sigmaU} is the unconditional
standard deviation used to calibrate the innovation scale; \code{sigmaE} is
the standard deviation of the log-variance innovation. The volatility process
uses the same grid choices as \code{discreteVAR}. An evenly spaced grid is
recommended when $\rho$ is large.

\subsection{Cox--Ingersoll--Ross process}

For the CIR diffusion
\begin{equation}
  \diff r_t=a(b-r_t)\diff t+\sigma\sqrt{r_t}\diff W_t,
\end{equation}
the transition density can be written
\begin{equation}
  f(r';r)=c\e^{-u-v}\left(\frac{v}{u}\right)^{q/2}
  I_q(2\sqrt{uv}),
\end{equation}
where
\begin{equation*}
  c=\frac{2a}{(1-\e^{-a\Delta})\sigma^2},\quad
  q=\frac{2ab}{\sigma^2}-1,\quad
  u=cr\e^{-a\Delta},\quad v=cr'.
\end{equation*}
The Feller restriction $\sigma^2<2ab$ is imposed. The conditional moments are
\begin{align*}
  \CE{r'|r}&=r\e^{-a\Delta}+b(1-\e^{-a\Delta}),\\
  \Cvar{r'|r}&=\frac{\sigma^2}{a}(1-\e^{-a\Delta})
  \left(r\e^{-a\Delta}+\frac b2\right).
\end{align*}

The interfaces are
\begin{lstlisting}[language=Matlab]
[P, X] = discretization.discreteCIR( ...
    a, b, sigma, Delta, N, Coverage, method);
\end{lstlisting}
\begin{lstlisting}[language=Python]
P, grid = discrete_cir(
    mean_reversion, long_run_mean, volatility, step,
    state_count, coverage=0.999, method="exponential",
)
\end{lstlisting}
Both evenly spaced and exponential grids are available. The exponential grid
is the default because the process is positive and its stationary distribution
can be concentrated near zero. The density itself is available through
\code{cirTransitionDensity} and \code{cir\_transition\_density}.

\section{Distribution and quadrature routines}

\subsection{Distribution from centered moments}

To construct an $N$-point distribution from centered moments, use
\begin{lstlisting}[language=Matlab]
[x, p] = discretization.momentMatchedDistribution( ...
    N, cMoments, q);
\end{lstlisting}
\begin{lstlisting}[language=Python]
grid, probability = moment_matched_distribution(
    state_count, centered_moments, prior=None
)
\end{lstlisting}
The first centered moment must be zero and the second must be positive. If the
prior is omitted, a Gaussian prior is used. The grid spans
$\sqrt{2N}$ standard deviations in each direction, as recommended by
\textcite{FarmerToda2017QE}.

\subsection{Gaussian-mixture quadrature}

The functions \code{gaussianMixtureQuadrature} and
\code{gaussian\_mixture\_quadrature} compute Gaussian-quadrature nodes and
weights for a Gaussian mixture using the Golub--Welsch construction. A
single-component input recovers ordinary Gaussian quadrature.

\subsection{Data-driven Gaussian quadrature}

The entry points are
\begin{itemize}
  \item \matlab{}: \code{dataDrivenGaussianQuadrature}; and
  \item \python{}: \code{data\_driven\_gaussian\_quadrature}.
\end{itemize}
Both functions construct quadrature nodes directly from one-dimensional data.
This method is separate from maximum-entropy moment matching and is based on
\textcite{Toda2021CE}.

An $N$-point rule requires moments through order $2N$. For positive or heavy-
tailed data, apply the method to logged observations when appropriate and
exponentiate the resulting nodes. This transformation reduces sensitivity to
high-order raw moments.

\section{Choosing grids and moments}
\label{sec:choices}

\begin{itemize}
  \item \textbf{Start with two moments.} Matching the conditional mean and
  variance usually provides a good balance between accuracy and feasibility.
  Higher moments require more grid support and can be infeasible near grid
  boundaries.
  \item \textbf{Use an evenly spaced grid for persistent processes.} A wider
  grid is often more valuable than a high-order quadrature rule when the
  conditional distribution shifts substantially across current states.
  \item \textbf{Use quadrature for weak persistence.} Gauss--Hermite or a
  model-specific quadrature prior can be effective when conditional means stay
  close to the center of the unconditional grid.
  \item \textbf{Avoid the VAR quantile grid for new work.} It is retained for
  replication and comparison.
  \item \textbf{Respect dimensional growth.} With $N$ points in each of $d$
  dimensions, the transition matrix contains $N^{2d}$ entries. Dense methods
  become impractical quickly even when each one-dimensional solve is cheap.
\end{itemize}

When a warning reports that $k$ moments could not be matched, the routine uses
the highest lower-order match that succeeded. The usual remedies are to widen
the grid, increase the number of points, or request fewer moments. Suppressing
the warning without checking the resulting moment errors is not recommended.

\section{Writing a custom discretization routine}
\label{sec:custom}

A one-dimensional custom routine has four steps:
\begin{enumerate}
  \item choose grid points and quadrature/integration weights;
  \item compute a strictly positive prior on that grid for each current state;
  \item define scaled conditional moment functions and targets; and
  \item call the maximum-entropy weight routine, checking the returned moment
  error and providing a lower-order fallback when needed.
\end{enumerate}

For example, the following code constructs weights on a fixed grid with mean
zero and variance one.

\begin{lstlisting}[language=Matlab]
D = linspace(-4, 4, 9);
T = @(x) [x; x.^2];
TBar = [0; 1];
q = normpdf(D);

[p, lambda, momentError] = ...
    discretization.maximumEntropyWeights(D, T, TBar, q);
\end{lstlisting}

The equivalent \python{} call is

\begin{lstlisting}[language=Python]
import numpy as np
from scipy.stats import norm
from discretization import maximum_entropy_weights

grid = np.linspace(-4, 4, 9)
moment_function = lambda x: np.vstack((x, x**2))
targets = np.array([0.0, 1.0])
prior = norm.pdf(grid)

probability, dual, moment_error = maximum_entropy_weights(
    grid, moment_function, targets, prior
)
\end{lstlisting}

Scaling moment functions by a representative grid magnitude generally improves
conditioning. Prior weights need not sum to one, but they must be nonnegative
with at least one positive entry. Very small positive floors are useful when a
density underflows in the tails.

\section{Numerical safeguards and verification}

The modernized code includes the following safeguards:
\begin{itemize}
  \item log-sum-exp evaluation of the dual objective;
  \item analytic gradient and Hessian calculations;
  \item finite-candidate fallback when an optimizer fails;
  \item convex-hull feasibility checks for selected moment targets;
  \item positive floors for prior weights;
  \item direct linear solves instead of explicit matrix inverses;
  \item a closed-form two-dimensional VAR variance rotation; and
  \item reduced temporary allocation in tensor-product transition matrices.
\end{itemize}

Automated tests cover row stochasticity, nonnegativity, default arguments,
historical compatibility wrappers, alternate grid methods, higher-dimensional
VAR structure, and reference parity between \matlab{} and \python{}. The
reference files were generated with \matlab{} R2024b. Representative
performance workloads are provided in the repository's \code{benchmarks}
directory; measured times are machine-specific.

\begingroup
\singlespacing

\section{Compatibility names}
\label{sec:migration}

Table \ref{tab:migration} maps historical \matlab{} calls to the canonical API.
The old entry points forward to the package implementation and issue migration
warnings. They remain available during modernization so that replication code
does not break unexpectedly.

\begin{table}[htbp]
\centering
\caption{MATLAB compatibility-name migration}
\label{tab:migration}
\footnotesize
\begin{tabularx}{\textwidth}{@{}>{\ttfamily}l >{\ttfamily\raggedright\arraybackslash}X@{}}
\toprule
\normalfont Historical name & \normalfont Canonical replacement \\
\midrule
discreteAR & discretization.discreteAR \\
discreteVAR & discretization.discreteVAR \\
discreteCIR & discretization.discreteCIR \\
discreteGMAR & \normalfont\matgmar \\
discreteSV & \normalfont\matsv \\
discreteNP & \normalfont\matdist \\
discreteApproximation & \normalfont\matmaxent \\
GaussianMixtureQuadrature & \normalfont\matgmq \\
NPGQ & \normalfont\matnpgq \\
CIRpdf & \normalfont\matcirpdf \\
\bottomrule
\end{tabularx}
\end{table}

\begingroup\raggedright
The earlier descriptive package names
\code{discreteARWith\allowbreak GaussianMixtureShocks},
\code{discreteARWith\allowbreak StochasticVolatility},
\code{discreteDistributionFrom\allowbreak Moments}, and
\code{nonparametricGaussian\allowbreak Quadrature} also remain as \matlab{}
compatibility wrappers.\par
\endgroup

\section{Citation and software responsibility}

Validate this research software by checking grid coverage, moment errors,
transition-row sums, and sensitivity to state placement. Report the language,
revision, routine, grid method, state count, and requested moment count.

Cite \textcite{FarmerToda2017QE} for the main method;
\textcite{TanakaToda2013EL} and \textcite{TanakaToda2015SINUM} for the
maximum-entropy foundations; and \textcite{Toda2021CE} for data-based
nonparametric Gaussian quadrature.

\footnotesize
\setlength{\bibitemsep}{0pt}
\printbibliography
\endgroup

\end{document}
